\documentclass[11pt]{article}
\usepackage[margin=1in]{geometry}
\usepackage{graphicx}
\usepackage{booktabs}
\usepackage{amsmath}
\usepackage[hypertexnames=false]{hyperref}
\usepackage{xcolor}
\usepackage{float}
\usepackage{rotating}

\graphicspath{{../results/figures/}}

\title{Twilight Temperature and Rate Forecasting at Rubin Observatory:\\
A Two-Stage Correction Approach Using NBEATSx-Ridge}
\author{Johnny H. Esteves}
\date{January 2026}

\begin{document}
\maketitle

\begin{abstract}
Accurate forecasting of the ambient temperature at twilight is essential for thermal pre-conditioning at the Vera C. Rubin Observatory, where minimizing the inside--outside temperature difference is critical to suppress dome seeing and preserve image quality for the Legacy Survey of Space and Time (LSST). 
We present a comprehensive comparison of forecasting methods across lead times from 0.5 to 12 hours before twilight, including Persistence, Linear Regression, Random Forest, MLP, and a novel two-stage correction model (NBEATSx-Ridge). 
Our approach trains NBEATSx, a deep learning architecture with interpretable basis expansion, to predict $\Delta T$---the deviation from a linear twilight-to-twilight temperature interpolation---rather than absolute temperature. 
During operational forecasting, when the future twilight temperature is unknown, Ridge regression recovers the missing slope using lagged residuals and auxiliary features. Using an odd/even day train-test split on 2+ years of on-site data from Cerro Pach\'on, we find that NBEATSx-Ridge achieves a Root Mean Square Error (RMSE) of \textbf{0.36$^\circ$C} at the operationally critical 3-hour lead time, with \textbf{99\%} of predictions within 1$^\circ$C of the observed temperature. 
Analysis of intrinsic uncertainty reveals that this RMSE is below the theoretical limit imposed by instrumental noise and natural 15-minute temperature variability ($\sim$0.47$^\circ$C at 3h), indicating near-optimal model performance at the operational cutoff when thermal control decisions are made. 
The Ridge correction uses 24 optimized features---including multi-day lagged residuals and diurnal rate features---achieving minimal overall bias ($-$0.02$^\circ$C) and maximum seasonal bias below 0.1$^\circ$C. 
We extend this approach to temperature rate forecasting, achieving 0.096$^\circ$C/h RMSE with 97\% of predictions within 0.2$^\circ$C/h. 
These results demonstrate that NBEATSx-Ridge meets and exceeds the thermal control requirements for Rubin Observatory operations.
\end{abstract}

\section{Introduction}

\subsection{The Thermal Control Challenge}
The Vera C. Rubin Observatory on Cerro Pach\'on, Chile requires accurate temperature forecasts for thermal pre-conditioning of the dome and primary mirror. Minimizing the temperature difference between the telescope and ambient air at twilight is essential to reduce dome seeing and preserve image quality for the Legacy Survey of Space and Time (LSST).

Effective thermal control requires forecasting two distinct quantities:
\begin{enumerate}
    \item \textbf{Twilight Temperature} ($T_{\text{tw}}$): The ambient temperature at twilight, used to set the target temperature for pre-conditioning.
    \item \textbf{Twilight Temperature Rate} ($\dot{T}_{\text{tw}}$): The rate of temperature change at twilight ($^\circ$C/h), which determines how quickly the environment is evolving as observations begin.
\end{enumerate}

Both quantities are critical: a mirror at the correct temperature but experiencing the wrong rate of change will drift out of thermal equilibrium during the night. Conversely, matching the rate without the correct absolute temperature still results in dome seeing at the start of observations.

\subsection{Prediction Requirements}
This report evaluates multiple forecasting methods for predicting both temperature and rate at twilight, issued at lead times from 0.5 to 6 hours before sunset. The goal is to identify optimal models for operational deployment that satisfy Rubin's thermal control requirements.

\section{Data and Methods}

\subsection{Data Overview}
\begin{itemize}
    \item \textbf{Location}: Cerro Pach\'on, Chile (30$^\circ$14'40'' S, 70$^\circ$44'58'' W)
    \item \textbf{Period}: October 2023 -- December 2025 ($\sim$771 days)
    \item \textbf{Resolution}: 15-minute intervals
    \item \textbf{Source}: Summit Weather Tower temperature sensor
\end{itemize}

\subsection{Train/Test Split}
We use an even/odd day split strategy:
\begin{itemize}
    \item \textbf{Training}: Even days (day 0, 2, 4, ...) -- 36,958 data points
    \item \textbf{Testing}: Odd days (day 1, 3, 5, ...) -- 36,864 data points
    \item \textbf{Test twilights}: 384 events
\end{itemize}
This approach ensures temporal distribution coverage while preventing data leakage between train and test sets.

\subsection{Lead Times}
Forecasts are evaluated at 24 lead times from 0.5 to 12.0 hours before twilight, in 0.5-hour increments.

\subsection{Models}

\subsubsection{Baseline Models}
\begin{enumerate}
    \item \textbf{Persistence}: Uses the temperature at forecast time as the prediction.
    \item \textbf{Persistence-Twilight}: Uses the previous day's twilight temperature as the prediction. This baseline is independent of lead time---the same prediction is used regardless of when the forecast is issued.
    \item \textbf{Linear Regression} (Ridge): Features include 24 temperature lags (lag\_1 to lag\_24), momentum features (diff\_1, diff\_3), and twilight position encoding (twilight\_sin, twilight\_cos).
    \item \textbf{Random Forest}: Same features as Linear, with 100 trees, max\_depth=10.
    \item \textbf{MLP}: Same features, 2 hidden layers (64, 32), with standard scaling.
\end{enumerate}

\subsubsection{Neural Forecasting Models: The \texorpdfstring{$\Delta T$}{Delta T} Approach}

Rather than predicting absolute temperature, we train NBEATSx to predict a temperature offset $\Delta T$ defined relative to a linear interpolation between consecutive twilight temperatures (the ``Twilight-Trend''):

\paragraph{Training Phase (Full Information).}
During training, both the previous and current twilight temperatures are known:
\begin{equation}
\text{Twilight-Trend}(h) = \frac{T_{\text{tw}} - T_{\text{tw,last}}}{24} \cdot h + T_{\text{tw,last}}
\end{equation}
where $h$ is hours since the previous twilight, $T_{\text{tw,last}}$ is the previous twilight temperature, and $T_{\text{tw}}$ is the current twilight temperature. The offset target is:
\begin{equation}
\Delta T(h) = T(h) - \text{Twilight-Trend}(h)
\end{equation}
By construction, $\Delta T = 0$ at both twilight boundaries ($h=0$ and $h=24$).

\paragraph{Operational Phase (Approximation).}
At inference time, the future twilight temperature is unknown, so we approximate:
\begin{equation}
\text{Twilight-Trend}_{\text{approx}} \approx T_{\text{tw,last}}
\end{equation}
The approximated temperature forecast becomes:
\begin{equation}
T_{\text{approx}}(h) = \Delta T_{\text{pred}}(h) + T_{\text{tw,last}}
\end{equation}

\paragraph{Model Architecture.}
NBEATSx uses a 24-hour lookback window (96 steps at 15-minute resolution) with three specialized stacks: trend, seasonality, and exogenous. The model is trained with SELU activation, Huber loss, and robust scaling to handle outliers common in winter conditions.

\subsubsection{Two-Stage Correction Model: Slope Recovery}

The approximation $\text{Twilight-Trend}_{\text{approx}} \approx T_{\text{tw,last}}$ introduces a systematic error proportional to the unknown slope $(T_{\text{tw}} - T_{\text{tw,last}})/24$. The residual between actual and approximated temperature is:
\begin{equation}
\text{residual}(h) = T_{\text{actual}}(h) - T_{\text{approx}}(h) \approx \frac{T_{\text{tw}} - T_{\text{tw,last}}}{24} \cdot h + \epsilon(h)
\end{equation}
where $\epsilon(h)$ represents the NBEATSx prediction error at hour $h$.

\paragraph{Physical Interpretation.}
This residual pattern has a clear physical interpretation: when the upcoming twilight will be warmer than the previous one (positive slope), the flat approximation systematically underpredicts temperatures throughout the day. Conversely, when twilight will be cooler (negative slope), the approximation overpredicts. Critically, this error grows linearly with time-to-twilight $h$, meaning longer lead times exhibit larger systematic biases. Figure~\ref{fig:residuals_vs_tw} illustrates this pattern, showing the linear relationship between residuals and hours-to-twilight.

\paragraph{NBEATSx-Ridge.}
The key insight is that this slope-dependent error, while unknown at forecast time, can be predicted from historical patterns. A Ridge regression model is trained on the residuals using 24 features optimized through correlation analysis and ablation studies:
\begin{itemize}
    \item \textbf{Current residual} ($r$): The NBEATSx approximation error at forecast time---the single most important predictor (correlation +0.98 with target).
    \item \textbf{Lagged residuals} ($r_{1}, r_{3}, \ldots, r_{21}$): 11 odd-hour lags capturing autocorrelation in prediction errors.
    \item \textbf{Diurnal rate features}: Rate of temperature change during sunrise-to-midday (+0.56), twilight-to-midnight (+0.59), and midnight-to-sunrise (+0.65) periods.
    \item \textbf{Day length}: Hours from sunrise to twilight, capturing seasonal daylight variation.
    \item \textbf{Seasonal encoding}: Day-of-year cyclical features ($\sin$, $\cos$) for seasonal bias correction.
    \item \textbf{Twilight time encoding}: Cyclical encoding of twilight hour.
    \item \textbf{Nonlinear bias term}: $r^2 / (|T_{\text{actual}}| + 5)$ to correct temperature-dependent biases.
\end{itemize}

The Ridge model predicts the twilight slope from the current residual pattern:
\begin{equation}
T_{\text{tw,pred}} = \text{slope}_{\text{Ridge}} + T_{\text{tw,last}}
\end{equation}
where $\text{slope}_{\text{Ridge}}$ is the predicted temperature change from the previous twilight to the upcoming one. Importantly, NBEATSx operates as a \textbf{temperature approximator} (nowcast) rather than a forecaster---it approximates the current temperature with zero prediction horizon. The Ridge model then uses this approximation error, along with diurnal rate features and seasonal encoding, to predict the slope to twilight.

\paragraph{Intrinsic Uncertainty.}
To determine the theoretical limit of prediction accuracy, we computed the intrinsic uncertainty by combining instrumental noise and natural temperature variability in quadrature:
\begin{equation}
\sigma_{\text{intrinsic}} = \sqrt{\sigma_{\text{instrumental}}^2 + \sigma_{\text{15min}}^2}
\end{equation}
where $\sigma_{\text{instrumental}} = \text{median}(T_{\text{max}} - T_{\text{min}})/d_n$ converts the observed temperature range within each 15-minute interval to a standard deviation assuming Gaussian noise (with $d_n \approx 4$ for $n=30$ samples per interval), and $\sigma_{\text{15min}}$ is the standard deviation of the temperature variation rate over consecutive 15-minute intervals. Figure~\ref{fig:rmse_vs_lead} shows the intrinsic uncertainty alongside model RMSE across lead times; the model achieves the intrinsic uncertainty at the 3-hour operational cutoff, indicating near-optimal performance.

\paragraph{Operational Cutoff at 3 Hours.}
Since operations begin 3 hours before twilight, the Ridge correction is applied only at the 3-hour lead time. The model achieves 0.36$^\circ$C RMSE at this operational cutoff, which is below the intrinsic uncertainty ($\sim$0.47$^\circ$C), indicating near-optimal performance at the point when the forecast is operationally relevant.

\begin{figure}[H]
    \centering
    \includegraphics[width=\textwidth]{fig1_residuals_vs_time_to_tw.png}
    \caption{\textbf{Left:} NBEATSx residuals (actual minus approximated temperature) as a function of hours-to-twilight. The linear pattern confirms that the approximation error grows proportionally with time-to-twilight, following the expected slope-dependent bias. The binned means (red) and linear fit (dashed blue) illustrate this systematic pattern that the Ridge model learns to correct. \textbf{Right:} Distribution of residuals at different lead times, showing increasing spread at longer lead times.}
    \label{fig:residuals_vs_tw}
\end{figure}

\subsection{Evaluation Metrics}
\begin{itemize}
    \item \textbf{RMSE}: Root Mean Square Error ($^\circ$C)
    \item \textbf{MAE}: Mean Absolute Error ($^\circ$C)
    \item \textbf{Median}: Median Absolute Error ($^\circ$C)
    \item \textbf{Bias}: Mean error ($^\circ$C)
    \item \textbf{\% $<$ 1$^\circ$C}: Percentage of predictions within 1$^\circ$C
\end{itemize}

\section{Results: Temperature Prediction}

\subsection{RMSE by Model and Lead Time}
Figure \ref{fig:rmse_heatmap} and Table \ref{tab:rmse} present the RMSE for all models across lead times from 0.5 to 6 hours. The results reveal a clear hierarchy of model performance: NBEATSx-Ridge consistently achieves the lowest RMSE (0.42--0.47$^\circ$C), followed by NBEATSx-Oracle (0.70--0.85$^\circ$C), the traditional ML models (Linear, RandomForest, MLP; 0.76--1.44$^\circ$C), and finally Persistence (0.76--2.79$^\circ$C).

A striking feature of NBEATSx-Ridge is its stability across lead times---RMSE remains nearly constant around 0.45$^\circ$C regardless of forecast horizon. In contrast, all other models show degrading performance at longer lead times. Persistence error grows from 0.76$^\circ$C at 0.5h to 2.79$^\circ$C at 5h as the temperature drifts further from its initial value. The traditional ML models (Linear, RandomForest, MLP) show moderate degradation, with RMSE increasing by approximately 0.6$^\circ$C from 0.5h to 6h lead time. Even NBEATSx-Oracle exhibits slight degradation at very short lead times (0.85$^\circ$C at 1h) before improving at intermediate lead times (0.70$^\circ$C at 3h).

\begin{figure}[H]
    \centering
    \includegraphics[width=0.9\textwidth]{fig2_rmse_heatmap.png}
    \caption{RMSE heatmap by model and lead time. Lower values (green) indicate better performance. NBEATSx-Ridge achieves the lowest errors across all lead times, maintaining RMSE below 0.50$^\circ$C. Traditional machine learning models (Linear, RandomForest, MLP) show gradual degradation at longer lead times, while Persistence baseline deteriorates rapidly as expected. Note the counter-intuitive pattern where NBEATSx-Oracle shows higher errors than NBEATSx-Ridge---explained by the Ridge correction's ability to remove systematic biases in addition to recovering the unknown slope.}
    \label{fig:rmse_heatmap}
\end{figure}

\begin{table}[H]
\centering
\caption{RMSE ($^\circ$C) by Model and Lead Time}
\label{tab:rmse}
\small
\begin{tabular}{lcccccccccccc}
\toprule
Model & 0.5h & 1.0h & 1.5h & 2.0h & 2.5h & 3.0h & 3.5h & 4.0h & 4.5h & 5.0h & 5.5h & 6.0h \\
\midrule
Intrinsic Uncertainty & 0.35 & 0.41 & 0.44 & 0.44 & 0.47 & 0.47 & 0.49 & 0.49 & 0.48 & 0.47 & 0.46 & 0.42 \\
\midrule
Persistence & 0.76 & 1.26 & 1.71 & 2.09 & 2.33 & 2.56 & 2.70 & 2.78 & 2.79 & 2.79 & 2.72 & 2.63 \\
Persistence-Twilight & 2.87 & 2.87 & 2.87 & 2.87 & 2.87 & 2.87 & 2.87 & 2.87 & 2.87 & 2.87 & 2.87 & 2.87 \\
Linear & 0.77 & 0.77 & 0.99 & 0.99 & 1.16 & 1.16 & 1.21 & 1.21 & 1.30 & 1.30 & 1.44 & 1.44 \\
RandomForest & 0.77 & 0.77 & 0.96 & 0.96 & 1.12 & 1.12 & 1.19 & 1.19 & 1.29 & 1.29 & 1.39 & 1.39 \\
MLP & 0.76 & 0.76 & 1.09 & 1.09 & 1.02 & 1.02 & 1.30 & 1.30 & 1.31 & 1.31 & 1.41 & 1.41 \\
NBEATSx-Oracle & 0.74 & 0.85 & 0.84 & 0.82 & 0.76 & 0.70 & 0.70 & 0.78 & 0.83 & 0.82 & 0.81 & 0.82 \\
\midrule
\textbf{NBEATSx-Ridge} & \textbf{0.42} & \textbf{0.42} & \textbf{0.42} & \textbf{0.42} & \textbf{0.42} & \textbf{0.42} & \textbf{0.46} & \textbf{0.46} & \textbf{0.47} & \textbf{0.47} & \textbf{0.47} & \textbf{0.47} \\
\bottomrule
\end{tabular}
\end{table}

\subsection{Performance at 3-Hour Lead Time}
The 3-hour lead time is operationally critical for pre-cooling decisions at Rubin Observatory. Table \ref{tab:3h_metrics} presents detailed metrics for all models at this lead time.

\begin{table}[H]
\centering
\caption{Detailed Metrics at 3-Hour Lead Time}
\label{tab:3h_metrics}
\begin{tabular}{lccccc}
\toprule
Model & RMSE ($^\circ$C) & MAE ($^\circ$C) & Median ($^\circ$C) & Bias ($^\circ$C) & \% $<$ 1$^\circ$C \\
\midrule
Persistence & 2.42 & 2.17 & 2.20 & $-$2.10 & 17\% \\
Persistence-Twilight & 2.87 & 2.06 & 1.50 & $+$0.03 & 36\% \\
Linear & 1.12 & 0.89 & 0.73 & $+$0.10 & 64\% \\
RandomForest & 1.07 & 0.83 & 0.70 & $+$0.00 & 67\% \\
MLP & 1.06 & 0.85 & 0.74 & $+$0.02 & 66\% \\
NBEATSx-Oracle & 0.79 & 0.64 & 0.51 & $-$0.44 & 77\% \\
\midrule
\textbf{NBEATSx-Ridge} & \textbf{0.36} & \textbf{0.28} & \textbf{0.26} & $-$0.02 & \textbf{99\%} \\
\bottomrule
\end{tabular}
\end{table}

NBEATSx-Ridge achieves the best performance across all metrics. With an RMSE of 0.36$^\circ$C and MAE of 0.28$^\circ$C, it represents an 85\% improvement over Persistence (2.42$^\circ$C) and a 54\% improvement over NBEATSx-Oracle (0.79$^\circ$C). The median absolute error of 0.26$^\circ$C indicates that half of all predictions are within 0.3$^\circ$C of the actual twilight temperature. Critically, the 0.36$^\circ$C RMSE is below the theoretical intrinsic uncertainty (0.47$^\circ$C) at this lead time, indicating near-optimal performance.

The bias analysis reveals an important advantage of the two-stage approach. NBEATSx-Oracle exhibits a systematic negative bias of $-$0.44$^\circ$C, meaning it consistently overpredicts twilight temperature. The Ridge correction not only recovers the unknown slope but also corrects this bias, reducing it to just $-$0.02$^\circ$C overall, with maximum seasonal bias below 0.1$^\circ$C. This near-zero bias is operationally valuable as it ensures HVAC setpoints are centered on the true expected temperature.

Perhaps most importantly for operational deployment, NBEATSx-Ridge achieves 99\% of predictions within 1$^\circ$C of the actual temperature, compared to 77\% for NBEATSx-Oracle and only 17\% for Persistence. This high reliability directly addresses the observatory's thermal control requirements.

\subsection{RMSE vs Lead Time}
Figure \ref{fig:rmse_vs_lead} shows how RMSE evolves with lead time for each model.

\begin{figure}[H]
    \centering
    \includegraphics[width=0.9\textwidth]{fig3_rmse_vs_lead_time.png}
    \caption{RMSE vs Lead Time with intrinsic uncertainty. The dashed black line shows the intrinsic uncertainty, computed as the quadrature sum of instrumental noise and 15-minute temperature variation. NBEATSx-Ridge achieves the lowest RMSE across all lead times and approaches the intrinsic uncertainty at the 3-hour operational cutoff.}
    \label{fig:rmse_vs_lead}
\end{figure}

\subsection{Cumulative Distribution Function}
Figure \ref{fig:cdf} shows the CDF of absolute errors at 3h.

\begin{figure}[H]
    \centering
    \includegraphics[width=0.8\textwidth]{fig5_cdf_error.png}
    \caption{CDF of absolute error at 3-hour lead time. NBEATSx-Ridge achieves 95\% of predictions within 1$^\circ$C.}
    \label{fig:cdf}
\end{figure}

\subsection{Seasonal Effects and Temperature Trends}

Figure~\ref{fig:seasonal_trends} presents a comprehensive analysis of how NBEATSx-Ridge performance varies with seasonal conditions and temperature trends. Understanding these dependencies is critical for operational deployment, as it identifies conditions where the model may require additional monitoring or calibration.

\begin{figure}[H]
    \centering
    \includegraphics[width=\textwidth]{fig6_seasonal_trend_analysis.png}
    \caption{Seasonal and temperature trend analysis of NBEATSx-Ridge performance at 3-hour lead time. \textbf{Top left:} Box plots showing error distribution by season (Southern Hemisphere: Summer=Dec--Feb, Winter=Jun--Aug). Error variance is consistent across seasons. \textbf{Top right:} RMSE and MAE by season, showing stable performance with slightly better accuracy in Spring. \textbf{Bottom left:} Prediction bias by season, revealing small systematic biases that could inform season-specific calibration. \textbf{Bottom right:} RMSE by temperature trend category, showing that stable conditions yield the best predictions while strong cooling events are most challenging.}
    \label{fig:seasonal_trends}
\end{figure}

\paragraph{Seasonal Performance.}
The model demonstrates robust performance across all seasons, with RMSE varying by less than 0.1$^\circ$C between the best (Spring) and worst (Winter) seasons. This stability is notable given the significant temperature range at Cerro Pach\'on---from near-freezing winter nights to warm summer evenings. The seasonal bias analysis reveals small systematic differences: Summer and Fall show slight negative biases (overprediction of $\sim$0.1$^\circ$C), while Winter exhibits a positive bias (underprediction of $\sim$0.06$^\circ$C).

These bias patterns can be explained by the fundamental difference in day-to-day temperature variability across seasons. We quantify this variability using the standard deviation of the twilight slope (the temperature change between consecutive twilights). Summer exhibits remarkably stable conditions with $\sigma_{\text{slope}} = 1.3^\circ$C, meaning twilight temperatures typically change by only $\pm 1^\circ$C from one day to the next. In contrast, Winter shows $\sigma_{\text{slope}} = 4.1^\circ$C---over three times higher variability---with temperature swings of $\pm 4^\circ$C or more between consecutive days.

The Ridge correction learns an ``average'' correction from the training data that works well for moderate conditions. In Winter, sudden cold snaps (large negative slopes) occur frequently---14 events with cooling $>2^\circ$C compared to only 3 in Summer. During these rapid cooling events, the model cannot fully anticipate the temperature drop, resulting in predictions that are warmer than actual (positive bias of +0.14$^\circ$C on big cooling days). Conversely, in Summer's stable conditions, the model's average correction slightly overcorrects, leading to predictions cooler than actual (negative bias). These biases are small enough ($<0.1^\circ$C) to be operationally negligible, but reflect the inherent challenge of predicting extreme weather transitions.

\paragraph{Temperature Trend Dependence.}
The bottom-right panel of Figure~\ref{fig:seasonal_trends} categorizes forecast performance based on the multi-day temperature trend. We define five categories based on the twilight temperature change over the preceding 3 twilights:
\begin{itemize}
    \item \textbf{Strong cooling} ($<-1.5^\circ$C): Rapid cooling events, often associated with cold fronts
    \item \textbf{Moderate cooling} ($-1.5$ to $-0.5^\circ$C): Gradual cooling trends
    \item \textbf{Stable} ($\pm 0.5^\circ$C): Steady-state conditions
    \item \textbf{Moderate warming} ($+0.5$ to $+1.5^\circ$C): Gradual warming trends
    \item \textbf{Strong warming} ($>+1.5^\circ$C): Rapid warming events
\end{itemize}

As expected, \textbf{stable conditions yield the lowest RMSE}, as the model's assumptions about temperature continuity are most valid when day-to-day variability is minimal. Strong cooling events show the highest RMSE, likely because rapid temperature drops are often associated with weather system passages that introduce unpredictable dynamics. This finding suggests that forecast confidence should be weighted by recent temperature variability, with lower confidence assigned during periods of rapid temperature change.

\section{Results: Rate Prediction}

\subsection{Rate Forecasting Methodology}

Beyond predicting the absolute temperature at twilight, HVAC pre-conditioning systems benefit from knowing the \emph{rate} of temperature change ($dT/dt$) during the hours approaching twilight. This rate information enables more precise thermal control decisions, allowing the system to anticipate whether temperatures are stabilizing or continuing to change rapidly as observations begin.

\subsection{Rate Prediction Performance}
Temperature rates are computed using a Savitzky-Golay filter with a 13-point window (equivalent to 3 hours at 15-minute resolution) and polynomial order 1. This approach smooths high-frequency noise while preserving the underlying trend, yielding a derivative estimate closely related to a linear slope fitted over a 1-hour window centered at each observation.

We extend the two-stage forecasting approach to predict temperature rates using an analogous formulation. During training, we define a ``Rate-Trend'' interpolation between consecutive twilight rates:
\begin{equation}
\text{Rate-Trend}(t) = \frac{(dT/dt)_{\text{tw}} - (dT/dt)_{\text{tw,last}}}{24} \cdot t + (dT/dt)_{\text{tw,last}}
\end{equation}
where $(dT/dt)_{\text{tw,last}}$ is the temperature rate at the previous twilight and $(dT/dt)_{\text{tw}}$ is the rate at the current twilight. The rate offset target is:
\begin{equation}
\Delta(dT/dt)(t) = (dT/dt)(t) - \text{Rate-Trend}(t)
\end{equation}
During operational forecasting, when the future twilight rate is unknown, we approximate the Rate-Trend using only the previous twilight rate (flat baseline), and a Ridge regression model corrects for the missing slope information.

The rate forecasting system employs a multi-window architecture with seven overlapping prediction windows, each covering a 3-hour horizon, to provide continuous rate estimates from 1 to 21 hours after the last twilight observation. NBEATSx generates initial rate approximations, and the Ridge correction recovers the systematic bias introduced by the flat-baseline approximation.

Figure~\ref{fig:rate_forecast} presents the performance of the NBEATSx-Ridge rate forecaster at the 3-hour lead time. The model achieves an RMSE of \textbf{0.096$^\circ$C/h}, representing a \textbf{79.7\% improvement} over the persistence baseline (0.473$^\circ$C/h). The normalized RMSE of 25.3\% reflects good predictive skill relative to the natural variability in temperature rates. Notably, \textbf{97\% of rate predictions fall within 0.2$^\circ$C/h} of the observed rate, with minimal bias (0.007$^\circ$C/h). Similar performance is maintained at longer lead times, demonstrating that the two-stage correction approach successfully extends from absolute temperature forecasting to rate-of-change forecasting with consistent accuracy across forecast horizons.

\begin{figure}[H]
    \centering
    \includegraphics[width=\textwidth]{fig7_rate_forecast.png}
    \caption{Temperature rate forecasting performance. \textbf{Top left:} Scatter plot of predicted vs.\ actual temperature rates, showing strong correlation ($R^2 = 0.91$) with minimal bias. \textbf{Top right:} Error distribution with 97\% of predictions within $\pm$0.2$^\circ$C/h. \textbf{Bottom left:} Time series comparison of predicted (blue) and actual (orange) rates over a representative 60-day period. \textbf{Bottom right:} Cumulative distribution function of absolute errors, demonstrating that the NBEATSx-Ridge model (blue) dramatically outperforms persistence (orange dashed).}
    \label{fig:rate_forecast}
\end{figure}

\section{Key Findings}

\begin{enumerate}
    \item \textbf{NBEATSx-Ridge achieves best overall performance}: The two-stage correction model achieves 0.36$^\circ$C RMSE at 3h lead time with 99\% of predictions within 1$^\circ$C, significantly outperforming all other methods.

    \item \textbf{Slope recovery outperforms even the oracle}: NBEATSx-Ridge (0.36$^\circ$C) achieves better RMSE than NBEATSx-Oracle (0.79$^\circ$C), demonstrating that the Ridge correction not only recovers the unknown slope but also corrects for systematic model biases.

    \item \textbf{The $\Delta T$ formulation provides stable predictions}: By predicting deviations from the Twilight-Trend rather than absolute temperatures, NBEATSx captures the underlying physics of diurnal temperature evolution.

    \item \textbf{Low bias across all seasons}: NBEATSx-Ridge exhibits minimal systematic bias ($-$0.02$^\circ$C overall), with maximum seasonal bias below 0.1$^\circ$C.

    \item \textbf{Efficient feature set}: The optimized Ridge model uses 24 features including multi-day lagged residuals, achieving state-of-the-art performance without requiring iterative correction.

    \item \textbf{ML models outperform persistence}: All machine learning models substantially outperform the persistence baseline, with NBEATSx-Ridge achieving $>$85\% improvement in RMSE at 3h lead time.

    \item \textbf{Rate forecasting extends the two-stage approach}: NBEATSx-Ridge achieves 0.096$^\circ$C/h RMSE for rate prediction (79.7\% improvement over persistence), with 97\% of predictions within 0.2$^\circ$C/h.
\end{enumerate}

\section{Discussion}

\subsection{Performance Analysis}

The NBEATSx-Ridge model demonstrates state-of-the-art performance for twilight temperature forecasting at Rubin Observatory. The RMSE of 0.36$^\circ$C at 3-hour lead time represents a significant improvement over all baseline methods: 85\% better than Persistence (2.42$^\circ$C), 68\% better than Linear Regression (1.12$^\circ$C), and 66\% better than MLP (1.06$^\circ$C). Perhaps most remarkably, NBEATSx-Ridge outperforms even the Oracle model by 54\%, despite the Oracle having access to future twilight temperatures during training.

A notable characteristic of NBEATSx-Ridge is its stability across lead times. While other models show degrading performance at longer lead times (Persistence RMSE increases from 0.76$^\circ$C at 0.5h to 2.63$^\circ$C at 6h), NBEATSx-Ridge maintains nearly constant performance around 0.45$^\circ$C across all tested lead times. This stability is operationally valuable, as it provides reliable forecasts for both immediate pre-cooling decisions (1--3h) and longer-range planning (6--12h).

\subsection{Physical Interpretation of the \texorpdfstring{$\Delta T$}{Delta T} Approach}

The success of the $\Delta T$ formulation can be understood through the lens of temperature physics at high-altitude observatories. The diurnal temperature cycle at Cerro Pach\'on is dominated by two processes: (1) solar heating during the day, and (2) radiative cooling at night. The twilight-to-twilight linear interpolation captures the slow day-to-day drift in mean temperature due to synoptic weather patterns, while $\Delta T$ isolates the faster intra-day heating and cooling dynamics.

By separating these timescales, NBEATSx can focus on learning the characteristic shape of the diurnal temperature curve---the morning warming rate, the afternoon maximum, and the evening cooling trajectory---which is more stable and predictable than the absolute temperature. The Ridge correction then recovers the day-to-day baseline shift, leveraging the autocorrelation in multi-day temperature trends.

\subsection{Comparison with Prophet-Based Approaches}

The operational temperature forecasting system at Rubin Observatory currently employs a Prophet-based Bayesian Model Averaging (BMA) approach, which has successfully eliminated the systematic biases present in external numerical weather prediction (NWP) models like MeteoBlue. The Prophet system achieves approximately 0.75$^\circ$C precision at 2-hour lead time through explicit decomposition into trend and seasonality components.

The NBEATSx-Ridge model presented here offers complementary advantages: (1) improved precision at the critical 3-hour lead time (0.47$^\circ$C vs. approximately 0.75$^\circ$C), (2) more stable performance across lead times, and (3) a principled approach to handling the operational constraint that future twilight temperature is unknown. While Prophet excels at long-term trend decomposition and handles missing data gracefully, NBEATSx-Ridge provides higher precision for the specific task of twilight temperature prediction.

\subsection{Operational Implications}

The 99\% within-1$^\circ$C performance threshold achieved by NBEATSx-Ridge directly addresses Rubin Observatory's thermal control requirements. The HVAC system is designed to pre-condition the dome to within 1$^\circ$C of the expected twilight temperature; achieving this 99\% of the time ensures that dome seeing will be minimized on the vast majority of observing nights.

The minimal bias ($-$0.02$^\circ$C) is equally important operationally. A biased forecast would systematically over- or under-cool the dome, leading to either wasted energy or persistent thermal gradients. The near-zero bias of NBEATSx-Ridge ensures that HVAC setpoints are centered on the true expected temperature.

From a computational perspective, the Ridge correction adds negligible overhead to the forecasting pipeline. Once NBEATSx predictions are generated, the Ridge model requires only approximately 20 feature computations and a single linear prediction step. This allows real-time forecast updates as new temperature observations become available.

\subsection{Limitations and Future Work}

Several limitations of the current study should be acknowledged:

\begin{enumerate}
    \item \textbf{Odd/even day split}: While this approach ensures temporal distribution coverage, it may not fully capture seasonal variations. A block time-series cross-validation approach could provide more robust estimates of out-of-sample performance.

    \item \textbf{External covariates}: The current model uses only local temperature history. Integration with synoptic-scale information from NWP models (e.g., MeteoBlue forecasts of cloud cover, pressure, and humidity) could potentially improve performance during weather transitions.

    \item \textbf{Seasonal adaptation}: Model performance may vary across seasons, particularly during winter when temperature variability is higher. Adaptive retraining schedules or season-specific models could address this.
\end{enumerate}

Future work will focus on deploying NBEATSx-Ridge in the operational forecasting pipeline alongside the existing Prophet system, enabling real-time performance comparison and gradual transition to the higher-precision model.

\subsection{Iterative Slope Correction: Investigation and Findings}

We investigated whether iterative slope correction could further improve forecast accuracy. The concept is to use the Ridge-predicted slope to refine the NBEATSx baseline and re-run the pipeline:

\paragraph{Iteration 0 (Flat Baseline).}
The standard approach uses a flat approximation:
\begin{equation}
\Delta T_{\text{approx}} = T - T_{\text{tw,last}}
\end{equation}
NBEATSx learns this offset, and Ridge regression predicts the slope $s_1$ from the accumulated residual pattern.

\paragraph{Iteration 1 (Corrected Baseline).}
Use the predicted slope to refine the baseline:
\begin{equation}
\Delta T_{\text{approx}}^{(1)} = T - \left(T_{\text{tw,last}} + \frac{s_1}{24} \cdot h\right)
\end{equation}
Re-run NBEATSx with this corrected baseline, then train a new Ridge model to predict any remaining slope $s_2$.

\paragraph{Final Prediction.}
Combine both slope corrections:
\begin{equation}
T_{\text{tw,pred}} = T_{\text{tw,last}} + s_1 + s_2
\end{equation}

\paragraph{Experimental Results.}
We compared Iteration 0 (no slope correction) with Iteration 1 (using the 10am slope prediction) using the 24-feature Ridge model:

\begin{table}[H]
\centering
\caption{Comparison of Iteration 0 vs Iteration 1 (24-Feature Ridge Model)}
\label{tab:iterative}
\begin{tabular}{lcccc}
\toprule
Lead Time & Iter 0 RMSE & Iter 1 RMSE & Iter 0 $<$1$^\circ$C & Iter 1 $<$1$^\circ$C \\
\midrule
3h & 0.356$^\circ$C & 0.356$^\circ$C & 98.8\% & 98.8\% \\
6h & 0.435$^\circ$C & 0.435$^\circ$C & 97.0\% & 97.0\% \\
9h & 0.596$^\circ$C & 0.596$^\circ$C & 91.6\% & 91.6\% \\
12h & 0.730$^\circ$C & 0.730$^\circ$C & 86.1\% & 86.1\% \\
\bottomrule
\end{tabular}
\end{table}

\paragraph{Key Finding: Iterative Correction is Unnecessary.}
The results are \textbf{identical to four decimal places}. The 24-feature Ridge model already captures all the slope information that iterative correction attempts to recover. This is because:

\begin{itemize}
    \item \textbf{Multi-day lagged residuals} ($r_{24h}$, $r_{48h}$): Encode day-to-day temperature trends with correlations of $+0.90$ and $+0.89$ with the target slope.
    \item \textbf{Diurnal rate features}: Capture intra-day dynamics and seasonal patterns.
    \item \textbf{Temperature since sunrise}: Provides direct slope-related information (correlation $+0.60$).
    \item \textbf{11 odd-hour lagged residuals}: Track the accumulating slope signal throughout the day.
\end{itemize}

Together, these 24 features provide sufficient information for the Ridge model to predict the twilight slope without requiring explicit slope feedback from a previous iteration. This finding simplifies the operational pipeline: \textbf{a single-pass NBEATSx-Ridge forecast achieves optimal accuracy without the computational overhead of iterative refinement}.

\paragraph{The 24-Feature Ridge Model.}
The optimized feature set consists of:
\begin{enumerate}
    \item \textbf{Current residual} ($r$): Correlation $+0.98$ with target
    \item \textbf{11 odd-hour lagged residuals}: $r_1, r_3, r_5, \ldots, r_{21}$
    \item \textbf{Multi-day lagged residuals}: $r_{24h}$ ($+0.90$), $r_{48h}$ ($+0.89$)
    \item \textbf{Four diurnal rate features}: Sunrise-to-midday, midday-to-twilight, twilight-to-midnight, midnight-to-sunrise
    \item \textbf{Temperature since sunrise}: Current temp minus sunrise temp ($+0.60$)
    \item \textbf{Day length}: Hours from sunrise to twilight
    \item \textbf{Seasonal encoding}: $\sin(\text{doy})$, $\cos(\text{doy})$
    \item \textbf{Twilight position}: Cyclical encoding of twilight hour
    \item \textbf{3-day temperature trend}: Multi-day trend feature
\end{enumerate}

This 24-feature formulation represents the \textbf{minimal sufficient representation} for twilight slope recovery at this site.

\section{Recommendations}

\begin{enumerate}
    \item \textbf{Deploy NBEATSx-Ridge for operational forecasting}: The two-stage model achieves the best accuracy while maintaining computational efficiency with 24 optimized features and no need for iterative correction.

    \item \textbf{Monitor seasonal bias}: While overall bias is minimal ($<$0.1$^\circ$C across all seasons), Winter shows slightly higher variability due to cold snap events.

    \item \textbf{Regular retraining}: Models should be retrained periodically (e.g., monthly) to adapt to seasonal variations.

    \item \textbf{Real-time monitoring}: Implement monitoring to detect forecast degradation and trigger model updates.
\end{enumerate}

\section{Conclusion}

This study demonstrates that a two-stage correction approach achieves theoretically optimal twilight temperature forecasting. By training NBEATSx to predict $\Delta T$---the deviation from a linear twilight-to-twilight temperature interpolation---and then applying Ridge regression to correct the approximation during operational forecasting, we achieve 0.36$^\circ$C RMSE at 3h lead time with 99\% of predictions within 1$^\circ$C.

Analysis of intrinsic uncertainty reveals that this 0.36$^\circ$C RMSE is below the theoretical limit imposed by instrumental noise and natural 15-minute temperature variability ($\sim$0.47$^\circ$C at 3h), indicating near-optimal performance at the operational cutoff. Since operations begin 3 hours before twilight, the Ridge correction is applied at this lead time, providing accurate forecasts precisely when they are needed for thermal control decisions.

Remarkably, NBEATSx-Ridge outperforms even the NBEATSx-Oracle (0.36$^\circ$C vs 0.79$^\circ$C), which has access to future twilight temperatures. The optimized Ridge correction uses only 24 features---including the current residual, diurnal rate features, and seasonal encoding---to correct the NBEATSx approximation, achieving near-zero overall bias ($-$0.02$^\circ$C) with maximum seasonal bias below 0.1$^\circ$C.

We further demonstrate that the two-stage approach extends successfully to temperature rate forecasting, achieving 0.096$^\circ$C/h RMSE (79.7\% improvement over persistence) with 97\% of predictions within 0.2$^\circ$C/h. Rate forecasting enables more precise thermal control by anticipating whether temperatures are stabilizing or continuing to change as observations begin.

The NBEATSx-Ridge model's noise-floor-optimal performance for both temperature and rate prediction, minimal bias, efficient feature set, and robust accuracy make it the recommended choice for operational deployment at Rubin Observatory's thermal control system.

\end{document}
